% 本科毕业论文科研初稿：使用真实CSV生成表图，学校模板待确认。
\IfFileExists{ctexbook.cls}{
  \documentclass[UTF8,12pt,oneside]{ctexbook}
}{
  \documentclass[12pt,oneside]{book}
  \usepackage{fontspec}
  \setmainfont{Noto Serif CJK SC}[AutoFakeSlant=.15]
  \setsansfont{Noto Sans CJK SC}
  \setmonofont{DejaVu Sans Mono}[Scale=.85]
  \XeTeXlinebreaklocale "zh"
  \XeTeXlinebreakskip=0pt plus 1pt
  \renewcommand{\contentsname}{目录}
  \renewcommand{\figurename}{图}
  \renewcommand{\tablename}{表}
  \renewcommand{\bibname}{参考文献}
}
\usepackage[a4paper,margin=2.5cm]{geometry}
\usepackage{amsmath,amssymb,amsthm,graphicx,booktabs,longtable}
\renewcommand{\proofname}{证明}
\usepackage[hidelinks]{hyperref}
\theoremstyle{definition}
\newtheorem{proposition}{命题}[chapter]
\newtheorem{remark}{注}[chapter]
\newcommand{\R}{\mathbb R}
\newcommand{\norm}[1]{\left\lVert#1\right\rVert}
\newcommand{\figurefile}[3]{\begin{figure}[htbp]\centering\includegraphics[width=.98\textwidth]{figures/#1.pdf}\caption{#2}\label{#3}\end{figure}}
\setlength{\parindent}{2em}
\setlength{\parskip}{.25em}
\linespread{1.25}
\emergencystretch=3em
\pagestyle{plain}
\title{基于自适应采样的局部随机神经网络求解椭圆型界面问题的数值研究\\[1em]\large 本科毕业论文科研初稿}
\author{作者、学号、学院及指导教师信息待填写}
\date{2026年10月13日}
\begin{document}
\frontmatter
\maketitle
\chapter*{摘要}
椭圆型界面问题具有分片系数及界面耦合条件，配置点在区域和界面上的分配会影响强形式离散系统的精度与稳定性。本文在既有局部随机神经网络思想基础上，独立实现固定隐藏参数的单隐层双子域方法，通过线性最小二乘确定输出系数。采用解析导数建立内部方程、外边界条件以及界面值和法向通量跳跃条件，并补充有限差分导数模式，对公开文献中的花形界面算例开展近似复现。

首先以正弦制造解验证Poisson方程五点差分和随机特征程序，差分观测收敛阶接近二阶。随后比较固定均匀点、固定随机点、仅PDE残差驱动采样及区域与界面联合指标采样，并开展预算、参数、消融和独立局部峰实验。所有随机比较至少采用五个固定种子，区分物理点预算与标量方程预算，在独立测试点报告解误差、界面误差及完整计算成本。

实验表明，在本文已测试的平滑直线界面问题与预设参数下，联合策略未取得相对均匀、随机基线的稳定解误差优势；增加界面资源的固定随机方案也能解释部分消融收益。局部峰问题中联合策略在部分对照改善，但多数情况下仍未优于均匀基线，且部分窄峰精度不足。尺度平衡、特征宽度、随机权重范围和SVD截断会显著改变误差及有效秩，解误差与物理通量误差可能具有不同排序。固定矩阵扰动实验进一步表明，系数敏感性与预测函数变化不可混为一谈。本文的贡献定位为已有方法的独立实现、可重复比较与有限范围稳定性分析，不主张提出具有普遍优势的原创采样算法。

\noindent\textbf{关键词：}椭圆型界面问题；局部随机神经网络；配置点；残差采样；最小二乘；数值稳定性
\chapter*{Abstract}
This thesis independently implements a two-subdomain randomized neural network collocation method for elliptic interface problems. Fixed hidden features are combined with analytic or finite-difference derivatives, and output coefficients are obtained by a linear least-squares solve. Manufactured solutions verify the PDE, boundary data, solution jumps, and coefficient-weighted normal flux jumps.

Uniform, random, PDE-residual, and joint interior-interface sampling are compared under explicit point or scalar-equation budgets. Independent tests, repeated seeds, ablations, parameter scans, localized peaks, and fixed-matrix perturbations are reported. The joint strategy does not consistently outperform the static baselines in the tested configurations. Narrow peaks expose approximation limitations, while scaling and rank truncation introduce accuracy trade-offs. A public flower-interface example is approximately reproduced with all implementation differences retained. The work emphasizes reproducible numerical comparison rather than a claim of algorithmic novelty or universal superiority.

\noindent\textbf{Keywords:} elliptic interface problem; randomized neural network; collocation; adaptive sampling; least squares; stability
\tableofcontents
\mainmatter
\chapter{绪论}
\section{研究背景与问题}
椭圆型偏微分方程常用于描述稳态扩散和势场。对于由不同材料组成的区域，扩散系数在材料界面发生跳跃，真解通常仅在每个子区域内光滑。即使解连续，梯度也可能不连续；在给定非零跳跃数据时，解本身亦可不连续。因此，单纯在整个区域内使用一个光滑近似并忽略界面条件，不能完整表达所要求解的问题。

传统差分或有限元方法通过网格及界面处理建立离散系统。随机神经网络方法则预先生成隐藏层参数，把近似写为固定特征的线性组合，只求输出系数。对于线性PDE，内部方程和线性边界、界面条件均可转化为线性约束。这为CPU上的可重复小规模研究提供了条件，但隐藏特征的相关性、配置点覆盖及约束尺度会影响求解结果。

本文研究的问题是：在固定随机特征和明确计算预算下，如何评价内部与界面配置点的选取及分配？残差较大的位置增加点数是否必然改善解误差？当系数对比度提高时，误差、矩阵条件性和计算成本是否具有一致的变化趋势？这些问题不能仅用训练损失回答，需要独立真解验证与合理静态基线。

\section{相关研究及贡献边界}
Sun、Dong和Wang将局部随机网络与间断Galerkin框架结合，提出多个格式并给出带假设的误差分析\cite{dgpreprint}。其中弱形式积分、面耦合及近似条件与本文强形式配置不同，不能将两者统称为同一离散格式。Li和Wang使用不同子域的随机网络求解界面问题，采用中心差分近似特征导数，并在随机点上叠置内部、界面和边界约束\cite{interfacepreprint}。其预印本总结明确把自适应配置点、初始化范围及数值分析列为后续方向；这为本文研究动机提供依据，却不保证本文方案在现阶段具有新颖性。

Dang、Wang和Jiang的AG-RaNN使用频率信息、残差和增长策略构造新增特征，并涉及域分裂\cite{ag}。本文阅读的是2025年的第三版预印本。其特征或网络架构适应与本文固定特征下的空间点适应不同；本文不直接套用其在特定正则性、参数及采样假设下的理论结论。Shang等结合单位分解、PCA与重叠Schwarz预条件来改善局部随机特征系统求解\cite{schwarzpreprint}，提示条件性不仅与点选取有关，也与特征冗余和求解器有关。

残差驱动采样本身已有明确先行研究。Wu等比较固定采样及RAD、RAR-D等方法\cite{rar}。较近的Residual-Christoffel方法依据算子作用后的特征空间几何构造密度、权重和白化\cite{christoffel}，其中“残差空间”并不等于当前近似解残差的幅值。本文的简单联合指标与这些技术的差异尚不足以支持原创算法声明。RNN-PG 2023论文正式书目信息已核实\cite{pg}，但全文尚未获得，因此本文不依据其未读推导建立算法或归纳作者结论。

\section{研究内容及组织}
本文的工作包括四个部分。第一，独立实现Poisson五点差分和随机特征配置，验证数学数据与导数。第二，构造零与非零跳跃的双子域制造解，研究加权最小二乘、尺度平衡与特征参数。第三，在共享特征和明确预算下比较四种采样及消融，保留无效、恶化和局部峰精度不足的结果。第四，给出标准有限维稳定性结论的自含推导，并通过扰动及花形界面近似复现讨论数值限制。

第二章介绍模型、制造解、差分和最小二乘；第三章推导局部随机特征离散系统；第四章说明采样及预算机制；第五章报告真实数值实验；第六章总结有限范围结论与未解决问题。附录给出补充模型与复现路径。

\section{研究范围与学术真实性}
本文不要求原创定理或重大算法创新。实验均为程序实际输出，原文报告数值与本项目数据分别列示。所有比较至少五种子，代表图不依据最好误差挑选。正式版和预印本存在题名、年份差异时分别记录；无法核实的全文内容不补写。本文尚未开展真实外部数据、三维、多层增长或完整Schwarz实现，也不作优于PINN或传统界面有限元的直接数值结论。

\chapter{预备知识与数学模型}
\section{椭圆型界面问题}
设单位正方形$\Omega=(0,1)^2$被$\Gamma=\{x=1/2,0<y<1\}$分为$\Omega_1$与$\Omega_2$，法向$n=(1,0)$从左指向右。系数$\beta_i>0$在子域内为常数。本文约定跳跃$[v]=v_2-v_1$，与参考预印本的相反约定须同时转换约束行和数据。
\begin{equation}\label{eq:interface}
\begin{cases}
-\beta_i\Delta u_i=f_i,&\text{在 }\Omega_i,\ i=1,2,\\
u=g_D,&\text{在 }\partial\Omega,\\
[u]=g_u,\quad[\beta\nabla u\cdot n]=g_q,&\text{在 }\Gamma.
\end{cases}
\end{equation}
本研究使用分片光滑制造解，把各项数据从同一函数严格导出。方程在每侧求导，不对跨界面不连续的总函数直接应用普通中心差分。

\section{Poisson制造解与差分验证}
令$u=\sin(\pi x)\sin(\pi y)$。由于$u_{xx}=u_{yy}=-\pi^2u$，有$f=2\pi^2u$及零Dirichlet边界。每向$n$个内部节点、$h=1/(n+1)$时五点格式为
\begin{equation}\label{eq:fdm}
\frac{4U_{ij}-U_{i-1,j}-U_{i+1,j}-U_{i,j-1}-U_{i,j+1}}{h^2}=f_{ij}.
\end{equation}
按$x$方向快序铺平，矩阵为$A=I\otimes T+T\otimes I$，其中$T=h^{-2}\operatorname{tridiag}(-1,2,-1)$。零边界不增加右端，非零边界则须把已知邻点贡献移至右端。

\begin{proposition}[正弦算例的离散误差]
对上述节点和制造解，五点差分解恰为
\[U=(2\pi^2/\lambda_h)u|_{\rm grid},\qquad\lambda_h=8h^{-2}\sin^2(\pi h/2).\]
相对离散$\ell^2$误差为$\pi^2h^2/12+O(h^4)$。
\end{proposition}
\begin{proof}
一维节点正弦是$T$的特征向量，特征值$4h^{-2}\sin^2(\pi h/2)$；二维张量积对应两个特征值之和。右端为$2\pi^2u|_{\rm grid}$，故解的倍数关系成立。由正弦平方的Taylor展开得
\[\lambda_h=2\pi^2-\pi^4h^2/6+O(h^4),\]
展开$2\pi^2/\lambda_h-1$即得结论。
\end{proof}
这是经典差分性质在特定真解上的初等推导，不是一般PDE或随机网络的新收敛定理。

\section{直线界面的分片制造解}
记$q=x(1-x)$、$p=q(x)(x-1/2)$，定义
\begin{align}
u_i&=\left(\sin\pi x+a_iq+\frac{t_i}{\beta_i}p\right)\sin\pi y,\label{eq:manufactured}\\
f_i&=\left[2\beta_i\pi^2\sin\pi x+\beta_i a_i(2+\pi^2q)-t_i(3-6x)+t_i\pi^2p\right]\sin\pi y.
\end{align}
因为$q''=-2$、$p''=3-6x$，式\eqref{eq:manufactured}满足内部方程。外边界为零，界面上$p=0,q=1/4,p'=1/4$，于是
\begin{equation}
g_u=\frac{a_2-a_1}{4}\sin\pi y,\qquad g_q=\frac{t_2-t_1}{4}\sin\pi y.
\end{equation}
连续模型取$a_1=a_2=0,t_1=t_2=1$，此时通量本身为$\sin\pi y/4$，并非恒零，能有效检验通量行。非零跳跃取$a_1=0,a_2=.4,t_1=1,t_2=1.2$，所以$g_u=.1\sin\pi y,g_q=.05\sin\pi y$。

\section{固定特征最小二乘与截断}
最小二乘、SVD和伪逆属于标准数值线性代数内容\cite{golub}。以下给出本文所需结论的自含推导，不将其列为原创定理。对实矩阵$B\in\R^{m\times n}$，求$\min_c\norm{Bc-d}_2^2$。SVD写为$B=Q\Sigma V^T$，保留$\sigma_j>\tau\sigma_1$的方向后，截断解为
\begin{equation}\label{eq:tsvd}
c_\tau=\sum_{\sigma_j>\tau\sigma_1}\frac{q_j^Td}{\sigma_j}v_j.
\end{equation}
满列秩时$\kappa_2(B)=\sigma_1/\sigma_n$。由正规矩阵特征值$\sigma_j^2$可知$\kappa_2(B^TB)=\kappa_2(B)^2$，这解释了为什么不直接用正规方程替代QR/SVD。本文使用SciPy的gelsd型求解，并单独记录奇异值诊断\cite{scipy}。

\begin{proposition}[固定矩阵右端扰动]
若$B$满列秩、$c=B^\dagger d$、$\widetilde c=B^\dagger(d+\delta d)$，则$\norm{\widetilde c-c}_2\le\norm{\delta d}_2/\sigma_n$。固定截断子空间时以最小保留奇异值替换$\sigma_n$。
\end{proposition}
\begin{proof}
差值为$B^\dagger\delta d$。由SVD，$\norm{B^\dagger}_2=1/\sigma_n$，应用算子范数不等式即得。固定截断空间的线性映射具有相同推导。若扰动沿最小保留左奇异向量，上界在精确算术中取到。
\end{proof}
这一标准线性代数结论不覆盖矩阵扰动或秩变化。若预测矩阵为$P$，函数变化是$P\delta c$，所以大的系数变化也不自动意味着大的函数变化。岭目标$\norm{Bc-d}^2+\alpha\norm c^2$对奇异方向使用$\sigma/(\sigma^2+\alpha)$，抑制放大但引入偏差，不能保证更小真解误差。

\chapter{局部随机神经网络离散方法}
\section{随机特征与物理导数}
每侧用$N_i$个固定特征$\phi_{ij}$，令$u_i^N=\sum_{j=1}^{N_i}c_{ij}\phi_{ij}$。隐藏权重与偏置随机生成后不优化。直线子域先映射到单位参考方形，左/右侧$\xi=(x-x_{\min})/.5,\eta=y$，特征为
\begin{equation}
\phi_j=\tanh\{w_{j1}(\xi-.5)+w_{j2}(\eta-.5)+b_j\}.
\end{equation}
若$t=\tanh z$，则$\sigma'=1-t^2$、$\sigma''=-2t(1-t^2)$。物理梯度带$x$方向因子2，物理Laplacian为$\sigma''(4w_{j1}^2+w_{j2}^2)$。忽略局部坐标尺度会产生模型错误，故程序以独立中心差分检查链式导数。

隐藏参数取指定均匀范围，不假设不同宽度的特征严格嵌套。相同种子改变$N_i$时，权重和偏置的随机调用顺序可能不同；宽度实验因此是同分布统计比较，不构成嵌套空间的单调性证明。

\section{分块矩阵组装}
令$c=(c_1^T,c_2^T)^T$。内部行分别为$(-\beta_1\Delta\Phi_1,0)$与$(0,-\beta_2\Delta\Phi_2)$，外边界仅激活所属侧的特征块。界面两类行是
\begin{equation}
A_u=(-\Phi_1,\Phi_2),\qquad
A_q=(-\beta_1\partial_n\Phi_1,\beta_2\partial_n\Phi_2).
\end{equation}
通量必须含各侧$\beta$，不能用导数连续代替通量条件。界面端点或角点可能造成多组约束重合，初期采用段中点或开区间随机点以避免重复。

各组样本数为$m_k$、平方权重为$\lambda_k$、固定物理尺度为$s_k$时，目标写为
\begin{equation}\label{eq:weighted}
\min_c\sum_k\frac{\lambda_k}{m_ks_k^2}\norm{A_kc-d_k}_2^2.
\end{equation}
叠置时矩阵和数据同时乘$\sqrt{\lambda_k/m_k}/s_k$。前期直线模型取边界与界面平方权重50；raw尺度全1，balanced内部尺度$\beta_i$、通量尺度$\beta_{\max}$。这些选择改变经验目标中的相对重要性，不是无损等价变换。高对比度时归一化通量小仍可对应较大物理通量误差，所以两种指标都需报告。

\section{求解与独立评估}
采用SVD型最小二乘避免显式正规方程，报告未截断条件数、保留条件数、有效秩、最小保留奇异值及系数范数。非零岭参数通过增广$[B;\sqrt\alpha I]$求解，数据矩阵秩与增广求解秩分开记录。

程序正常且指标有限仅说明运行完成，不代表精度验收或理论收敛。解误差在独立点上计算；均匀等权点的相对离散误差为$\sqrt{\sum|u_N-u|^2/\sum|u|^2}$，并报告最大误差。PDE残差使用解析算子在独立点计算。训练损失可能很小而边界、界面或未采区域误差仍大，有限点零残差不能保证连续PDE被满足。

\section{有限差分导数模式与原文区别}
对每个子域的平滑特征采用
\begin{align}
\partial_x\phi&\approx\frac{\phi(x+h_1,y)-\phi(x-h_1,y)}{2h_1},\\
\partial_{xx}\phi&\approx\frac{\phi(x+h_2,y)-2\phi(x,y)+\phi(x-h_2,y)}{h_2^2}.
\end{align}
即使点靠近界面，也只对同一网络特征延拓求差分，不切换到另一侧网络。对足够光滑函数，两种中心格式的截断误差为$O(h^2)$；二阶差分中浮点相减误差会除以$h_2^2$放大，因此步长并非越小越好。

花形近似复现按原文使用物理坐标$z=w\cdot x+b$及行权重$\gamma=50$\cite{interfacepreprint}，没有沿用直线模型的参考坐标、样本数归一化或平方权重50。其初始化及积分细节见第五章。本文强形式实现不是DG或PG方法，也不将解析导数版本称为原文FD严格复现。

\chapter{配置点策略与预算分配}
\section{基线及公平比较}
固定隐藏特征后，配置点决定哪些残差和约束被经验目标看到。本文比较A固定完整矩形网格、B固定随机点、C内部PDE残差增量点及D区域与界面联合指标。A/B一次求解；C/D先以共同初始点求解，再在同一候选池中新增点并重求输出系数，隐藏参数保持不变。

在同一随机种子内，特征流、点组流、候选流和抽样流分别管理。B的最终随机点包含C/D初始点前缀；A使用规则网格，不宣称与随机初始集相同。所有方法使用同一外边界点、固定权重和最终测试点。五种子为11、23、37、51、71，所有运行及失败状态均保存，不能只取最好一次。

界面一个物理点提供值与通量两行，故
\begin{equation}
m_{\rm eq}=m_{I,1}+m_{I,2}+2m_\Gamma+m_B.
\end{equation}
固定物理点预算时D可改变方程数；固定方程预算时D可改变物理点数。本文分别试验并同时报告两种计数，避免把新增界面方程的收益误归因于选点位置。

\section{残差指标和候选抽样}
给定当前解，内部残差$r_i=-\beta_i\Delta u_i^N-f_i$，界面残差$r_u=[u^N]-g_u$、$r_q=[\beta\partial_nu^N]-g_q$。在固定尺度下定义
\begin{equation}\label{eq:indicators}
\eta_i=\operatorname{RMS}(r_i/\beta_i),\qquad
\eta_\Gamma=\operatorname{RMS}\sqrt{r_u^2+(r_q/\beta_{\max})^2}.
\end{equation}
这些是采样指标而非已证明的后验解误差估计器。高对比度归一化可能弱化某些物理量的重要性，故实验仍报告未归一化通量。

C按残差分布选择内部新增点，但组预算和界面新增数固定；D根据三个指标调整预算，再按各组位置残差选点。可用候选点上的概率为
\begin{equation}
p_\ell=\frac{.2}{M}+.8\frac{|r_\ell|}{\sum_j|r_j|},
\end{equation}
界面用式\eqref{eq:indicators}中相应点指标代替$|r_\ell|$。全零时退回均匀概率，抽样无放回。这种探索混合避免只集中在少数候选点，但不保证连续区域覆盖或最佳配置。

非均匀点改变经验目标所偏重的位置。本文没有逆密度权重，因此不声称近似同一个固定均匀积分；连续期望可写为$\int p(x)r(x)^2\,dx$，与均匀残差范数不同。残差概率及增量/替换采样思想已有先行研究\cite{rar}；本文只研究一种简单联合预算机制的实际表现。

\section{最低配额及整数预算}
每轮最低新增点$(16,16,8)$，成本向量在点预算下为$(1,1,1)$，在方程预算下为$(1,1,2)$。扣除最低配额后，剩余成本按$\eta_k+\varepsilon$比例分配，除以各组点成本得到实数配额，先取整，再以剩余成本和最大欠配额补齐。成本1组保证此实验中预算能用尽。

初始内部点128/侧、界面32。静态增量为$(80,80,32)$，对应每轮192点或224标量行。主平滑比较三轮，A/B/C最终内部368/侧、界面128；峰值比较增加六轮设置，并固定最终总行1232或1904，外边界240行不变。

完整流程为：\begin{enumerate}
\item 固定模型、特征、边界、预算、候选池及随机流，保存配置。
\item 使用初始点求解；在候选池计算当前残差，不使用真解误差选点。
\item 计算组指标及整数配额，按探索混合概率新增未使用点。
\item 重新组装并求输出系数，重复预设轮数，不按最终测试误差挑轮次。
\item 在从未参与选点的网格评估全部指标，保留配额轨迹、系数及完整时间。
\end{enumerate}

\section{消融、候选偏差及复杂度}
去PDE或去界面指标时，对应组指标不参与动态预算，位置抽样退回均匀，最低配额和总预算仍不变。这个消融会同时改变配额与位置，因此后续固定比例扫描是必要的混杂因素对照，不能把消融变好直接解释为某指标无价值。

原D下一轮组RMS仍包含已经选入训练的候选点，可能乐观。D\_unused仅排除这些点后计算组RMS；位置选择仍无放回。该未用池也受历史选择影响，不是独立同分布最终测试集。本文没有据此使用iid采样的理论保证。

总特征$K$、矩阵行$m$、候选点$M$时，固定维度下组装约$O(mK)$，稠密SVD在$m\ge K$时主要为$O(mK^2)$，候选评估约$O(MK)$。$T$轮增量方案包含$T+1$次求解及$T$次候选评估，必须与静态的一次求解比较完整成本。重复奇异值诊断也纳入本项目计时；不能只展示最终一次求解的时间。

\chapter{数值实验与结果分析}
\section{环境、指标与数据来源}
实际运行环境为Linux CPU、Python 3.12.14、NumPy 2.2.6、SciPy 1.15.3及Matplotlib 3.10.3，BLAS线程设为1。程序设计为Windows Python 3.12可运行，但尚未完成Windows实机验证。无需GPU或真实外部数据。直线模型多数使用五种子，花形近似复现主比较使用十种子；表中的均值覆盖全部指定运行，样本标准差采用分母$n-1$。

直线模型测试每侧83乘83中点及257界面点；峰值细测试每侧161乘161、界面521点，另用81乘81检查积分敏感性。全局Poisson随机特征使用81乘81内部点。FDM误差在其节点上，残差为离散系统相对残差；其他PDE残差为独立点上的连续解析算子残差，两者不能直接作为同指标比较。时间区分单次组装/求解和完整自适应过程，不含最终验证/绘图。

所有表格由原始CSV生成，图由CSV重绘。主要来源为实验01、02、03、04、06、07、08及09的两组复现记录，均为2026年10月13日保存。完整路径见附录；原文数据另标注，不混为本项目输出。

\section{Poisson基线正确性}
表\ref{tab:fdm}和图\ref{fig:fdm}显示四级网格的误差与观测阶。相邻步长减半后误差约降四倍，与第二章的正弦特征值推导一致。单元测试还核查矩阵符号、制造右端、tanh导数和秩亏最小范数解。
\begin{table}[htbp]\centering\caption{Poisson五点差分真实结果}\label{tab:fdm}\begin{tabular}{lll}\toprule
每向内部点 & 相对离散误差 & 观测阶 \\ \midrule
15 & $3.219\times10^{-3}$ & -- \\
31 & $8.036\times10^{-4}$ & 2.0021 \\
63 & $2.008\times10^{-4}$ & 2.0005 \\
127 & $5.020\times10^{-5}$ & 2.0001 \\
\bottomrule\end{tabular}
\end{table}
\figurefile{fdm_convergence}{正弦Poisson五点差分收敛；数据来自exp01原始CSV。}{fig:fdm}
基础全局RaNN每种特征数五次：64特征相对误差均值约$8.64\times10^{-3}$，128特征约$5.99\times10^{-5}$，后者样本标准差约$3.08\times10^{-5}$。与此同时128特征加权矩阵平均条件数约$2.48\times10^9$。精度改善并不意味着条件性改善。全局网络与双子域网络总特征数、坐标尺度、配置点及评估方式不同，本文不据此宣称公平效率优越。

\section{直线界面与尺度取舍}
150次运行包含连续与非零跳跃、七组正反向系数组合、raw/balanced及Poisson极限。无运行异常，平衡组通过预设粗误差门槛；这一门槛不覆盖后续窄峰。表\ref{tab:interface}为非零跳跃、每侧128特征、scale2的五次均值。
\begin{table}[htbp]\centering\caption{高对比度下解与物理通量的尺度取舍}\label{tab:interface}\begin{tabular}{llll}\toprule
$(\beta_1,\beta_2)$ & 尺度 & 相对$L^2$均值 & 通量RMS均值 \\ \midrule
(1000,1) & raw & $4.579\times10^{-5}$ & $1.763\times10^{-4}$ \\
(1000,1) & balanced & $1.391\times10^{-7}$ & $2.756\times10^{-4}$ \\
(1,1000) & raw & $1.094\times10^{-4}$ & $1.711\times10^{-4}$ \\
(1,1000) & balanced & $1.466\times10^{-7}$ & $4.453\times10^{-4}$ \\
\bottomrule\end{tabular}
\end{table}
\figurefile{interface_contrast}{非零跳跃、每侧128特征的系数对比度实验；物理通量未除以系数。}{fig:interface}
平衡后解误差明显降低，但物理通量均值不一定降低。raw最高条件数约$1.69\times10^{15}$、最低有效秩240/256；balanced仍可达到约$2.11\times10^{12}$、最低255/256。截断SVD使近秩亏系统可求解，但不能据正常运行声称病态已消除。

\section{采样主比较与消融}
180次主比较固定每侧64特征、scale2、三轮；每种预算覆盖三个系数组合和五种子。表\ref{tab:sampling}展示点预算、$\beta=(1000,1)$的全部六种方法，最终物理总点1104，D方程数可变。
\begin{table}[htbp]\centering\caption{等物理点预算下采样与消融，五种子均值}\label{tab:sampling}\begin{tabular}{llll}\toprule
方法 & 相对$L^2$均值 & 完整时间/秒 & 界面点均值 \\ \midrule
A & $1.271\times10^{-4}$ & 0.0133 & 128.0 \\
B & $1.200\times10^{-4}$ & 0.0137 & 128.0 \\
C & $1.555\times10^{-4}$ & 0.0763 & 128.0 \\
D & $1.591\times10^{-4}$ & 0.0767 & 68.0 \\
去PDE指标 & $1.103\times10^{-4}$ & 0.0829 & 512.0 \\
去界面指标 & $1.561\times10^{-4}$ & 0.0743 & 56.0 \\
\bottomrule\end{tabular}
\end{table}
\figurefile{sampling}{点预算下全部种子的解误差，横线为均值；数据来自exp03。}{fig:sampling}
在每种预算的3设置乘5种子配对中，D对A/B解误差胜出均为0/15，D/B误差比例中位数约1.264和1.256。设置之间相关，不能把15个配对当作15个独立种子做显著性推断。D对C仅部分改善，仍没有稳定胜出。去PDE指标消融给界面更多点而较好，但后续相同最终组配额的固定随机在六组均值中更好且更快，未支持把收益简单归因于残差选点。

这说明比较基线必须合理：若只选界面明显欠采样的静态比例，可能高估动态分配的价值。另一方面，D默认指标偏向内部，界面配额较低，也可能影响解误差。本研究把该解释保留为数值猜想，而非因果证明。

\section{参数、秩与误差饱和}
285次探索性后续实验包含固定比例及单因素扫描。表\ref{tab:width}为$\beta=(1,1000)$宽度对照，其他设置固定。
\begin{table}[htbp]\centering\caption{特征宽度敏感性，数据秩上限为总特征数}\label{tab:width}\begin{tabular}{llll}\toprule
每侧特征 & 相对$L^2$均值 & 原始条件数均值 & 数据秩范围 \\ \midrule
32 & $1.138\times10^{-2}$ & $1.823\times10^{5}$ & 64--64 \\
64 & $1.085\times10^{-4}$ & $1.516\times10^{8}$ & 128--128 \\
128 & $1.236\times10^{-7}$ & $9.579\times10^{11}$ & 255--256 \\
256 & $2.581\times10^{-10}$ & $1.076\times10^{17}$ & 404--417 \\
\bottomrule\end{tabular}
\end{table}
\figurefile{width}{三个系数组合下的宽度、物理通量与条件数；数据来自exp04。}{fig:width}
当前平滑模型增宽改善误差，但条件数升至约$10^{17}$，256/侧时有效秩仅404至417，不能以总特征数代表独立维度。scale从.5到8的影响显著，小范围可适合平滑函数却加重病态，大范围可能使近似质量变差。界面权重、岭和点数也改变排序：强岭增加偏差；点数216至1728的解误差不单调，物理通量却可改善。这些扫描属于看到主比较后登记的探索性研究，不把其中最好参数改写为事先最优基线。

\section{独立局部峰与多预算}
在原分片制造解两侧加$h=16A x(1-x)y(1-y)\exp(-((x-c_x)^2+(y-c_y)^2)/w^2)$，取$A=.75,(c_x,c_y)=(.46,.55)$，两峰宽$.12,.06$。值跳跃不变，但通量增加$(\beta_2-\beta_1)h_x$，须重新制造数据。400次比较采用每侧128特征、scale2/4、两个行预算及五方法五种子。两个测试网格误差相对变化最大约0.49\%，未进入选点。
\figurefile{peaks}{scale4下独立峰的两个标量行预算比较；所有预设模型与方法均保留。}{fig:peaks}
80描述性配对中，D对A、B、C误差更低分别6、36、49次；对应比值中位数约1.293、1.095、.963。D\_unused没有稳定逆转结果。部分窄峰相对误差大于1，最大约1.493，无运行异常并不代表精度验收。固定特征表达力和采样目标变化都可能影响结果，不能只增加点数就承诺解决。

主峰值前段与谱任务并发，严格成本讨论使用另行25次串行重测。表\ref{tab:peaks}为固定窄峰、$\beta=(1000,1)$、scale4、六轮及1904方程，系数已与主case校验一致。D在该设置比C精度更好，但仍不如A且更贵。
\begin{table}[htbp]\centering\caption{固定峰值场景的串行完整成本与误差}\label{tab:peaks}\begin{tabular}{lll}\toprule
方法 & 相对$L^2$均值 & 串行完整时间/秒 \\ \midrule
A & $1.948\times10^{-1}$ & 0.0481 \\
B & $3.839\times10^{-1}$ & 0.0450 \\
C & $3.524\times10^{-1}$ & 0.4037 \\
D & $2.845\times10^{-1}$ & 0.4097 \\
D\_unused & $2.742\times10^{-1}$ & 0.4122 \\
\bottomrule\end{tabular}
\end{table}

\section{固定矩阵扰动观察}
15组矩阵谱、三阈值和两个扰动方向构成90组观察。右端相对扰动固定$10^{-8}$；最小保留左奇异方向的放大与$1/\sigma_{\min,\rm retained}$在浮点误差内一致，图\ref{fig:perturb}展示系数和函数变化的差别。
\figurefile{perturbation}{固定TSVD空间下的右端扰动；理论界仅针对固定矩阵和保留空间。}{fig:perturb}
最大系数相对变化约275.6，而预测函数相对变化最大约$1.71\times10^{-8}$。冗余特征可使系数敏感而函数变化较小，但该数值还可能接近大系数抵消的浮点求值地板，未以高精度算术复核。不能由此断言病态无害，也没有考察矩阵扰动、阈值跨秩或所有物理数据误差。

\section{花形界面FD近似复现}
选Li--Wang预印本Example4.1\cite{interfacepreprint}，$\Omega=(-1,1)^2$，中心$.02\sqrt5$，极径$r=.4+.2\sin20\theta$。图1核查花内为$\Omega_1$。两侧真解为$e^{xy}/\beta_1$及$\sin x\sin y/\beta_2$，$\beta=(1,10)$；法向由曲线切向旋转并归一化，不能只用径向单位向量。

按最贴合文字的解释，两网络权重范围均为1.6、偏置范围.7。与原文相同的设置包括单隐层tanh、3:1:1点数比例、$h_1=10^{-6},h_2=5\times10^{-4}$、行权重50、GL20及十次平均。作者种子、dtype、最小二乘阈值、确切界面测度未获实现；本项目float64/gelsd/$\tau=10^{-12}$、主角度均匀并补弧长对照。因此\textbf{只称近似复现}。
\begin{table}[htbp]\centering\small\caption{原文表1与本项目近似复现，十种子均值}\label{tab:reproduction}\resizebox{\textwidth}{!}{\begin{tabular}{lllll}\toprule
$m/N$ & 导数 & 原文表1 & 本项目GL20 & 本项目GL100 \\ \midrule
80/1000 & fd & $1.180\times10^{-3}$ & $3.169\times10^{-5}$ & $3.236\times10^{-5}$ \\
80/1000 & analytic & $1.180\times10^{-3}$ & $3.169\times10^{-5}$ & $3.236\times10^{-5}$ \\
160/1000 & fd & $9.200\times10^{-6}$ & $1.059\times10^{-7}$ & $1.080\times10^{-7}$ \\
160/1000 & analytic & $9.200\times10^{-6}$ & $1.059\times10^{-7}$ & $1.080\times10^{-7}$ \\
320/5000 & fd & $3.170\times10^{-8}$ & $3.057\times10^{-10}$ & $3.076\times10^{-10}$ \\
320/5000 & analytic & $3.170\times10^{-8}$ & $2.950\times10^{-10}$ & $2.961\times10^{-10}$ \\
\bottomrule\end{tabular}
}\end{table}
\figurefile{reproduction}{最贴合文字初始化下的FD与解析导数对照；原文曲线不是本项目数据。}{fig:reproduction}
表\ref{tab:reproduction}的误差与原文明显不同。更低数值不能当作算法改进，差异可能涉及初始化、dtype、截断与采样，尚无法唯一归因。最初按子域解释范围的110次保留为独立对照，核对文字后追加60次权重/偏置解释，旧数据不改写。原文通量印刷矩阵疑似漏$\beta$，本项目按物理模型保留，正式版仍待核对。

在子域范围解释变体中，五种子的FD步长结果见表\ref{tab:steps}。过大步长有截断偏差，过小步长出现舍入放大；不能把此变体当作作者结论。GL20与GL100在极小误差下可有明显相对变化，均保留而不选择较好的积分值。
\begin{table}[htbp]\centering\caption{FD步长独立诊断，$m=160,N=1000$}\label{tab:steps}\begin{tabular}{lll}\toprule
二阶差分步长 & 相对$L^2$均值 & 解析PDE残差RMS \\ \midrule
$1.000\times10^{-2}$ & $2.635\times10^{-8}$ & $4.659\times10^{-6}$ \\
$5.000\times10^{-4}$ & $5.976\times10^{-10}$ & $2.213\times10^{-7}$ \\
$1.000\times10^{-5}$ & $1.262\times10^{-7}$ & $1.768\times10^{-5}$ \\
$1.000\times10^{-6}$ & $8.669\times10^{-6}$ & $8.555\times10^{-4}$ \\
\bottomrule\end{tabular}
\end{table}

\section{实验结论的范围}
已有数据支持对采样、尺度、表达力和求解截断进行系统讨论，但不支持普适自适应优势。平滑、峰值及花形模型具有不同的近似难度；对一个算例表现好的范围不应直接迁移。还未做传统界面FDM/FEM或PINN同预算对照、完整多层增长、复杂物理或矩阵扰动。初稿中所有结论限于已登记模型与参数，后续应以独立模型和理论补充，而非对旧测试集无限调参。

\chapter{总结与展望}
\section{主要工作与结论}
本文建立了CPU可运行、配置和原始结果可追溯的随机特征科研工程，验证Poisson五点差分、双子域解析导数及非零界面跳跃，并实现固定均匀、固定随机、PDE残差和联合指标采样。数学分析包括离散矩阵、局部尺度、加权最小二乘、标准右端扰动界及岭滤波解释。花形界面FD近似复现补充了与原文实现细节不同的诊断。

第一，合理静态基线不可缺少。当前平滑模型的联合策略没有稳定胜过均匀或随机点，固定更多界面点也可取得消融中的部分收益，不能只依据一个较差固定比例宣称方法改进。

第二，采样不能脱离表达力和目标解释。独立窄峰暴露了固定特征精度不足；联合策略在部分对照改善而未取得普遍优势。非均匀经验目标、配额和位置选择的效应需要区分，单纯增加点数的误差也不必单调。

第三，条件数、系数敏感性和物理解误差不是同一指标。尺度平衡可改善解却恶化物理通量；增宽可降误差同时增加病态；SVD截断会改变有效近似空间。固定矩阵的系数放大界有数值支持，但小函数变幅不足以证明一般稳定性。

第四，复现必须保留不一致。本文确认了原文花形数据和FD设置，却未取得足以逐值严格重现的全部细节，故明确标为近似复现。与原表更低或更高的误差都不自动构成新贡献。

\section{局限及后续方向}
研究仍局限于少量二维制造解、固定单隐层特征和五至十种子；正式版公式、RNN-PG全文和完整新颖性核查有缺口。实验没有建立随机自适应的全PDE误差理论，也没有完成学校官方模板确认。窄峰结果提示未来可研究更多局部子域、特征增长或算子感知采样，但这些已有相关工作\cite{ag,christoffel,schwarzpreprint}，应先明确差异，且把额外表达力和成本公平地给予对照。

后续优先处理独立监测池、尺度与重要性权重的关系、矩阵扰动及高精度抵消核查，并在新的模型上验证从已有结果形成的假设。论文的合理定位仍为既有方法的独立实现、可重复数值比较和稳定性分析；是否形成算法改进，应由更强证据决定。

\appendix
\chapter{制造解、程序与复现说明}
\section{局部峰导数}
令$q(s)=s(1-s)$、$\xi=x-c_x,\eta=y-c_y,a=w^{-2}$，则
\begin{align}
h_x&=16A e^{-a(\xi^2+\eta^2)}[q'(x)-2a\xi q(x)]q(y),\\
h_{xx}&=16A e^{-a(\xi^2+\eta^2)}[-2-4a\xi q'(x)+(4a^2\xi^2-2a)q(x)]q(y).
\end{align}
$h_{yy}$交换$x,y$即可；$f_i^{\rm peak}=f_i-\beta_i\Delta h$。两侧共同$h$在值跳跃中相消，通量却因$\beta$不同而改变。

花形曲线$X=c+r(\theta)(\cos\theta,\sin\theta)$的切向为$(r'\cos\theta-r\sin\theta,r'\sin\theta+r\cos\theta)$。其外法向为切向顺时针旋转后的单位向量，分母$\sqrt{r^2+(r')^2}$。内部强迫为$-(x^2+y^2)e^{xy}$，外部强迫为$2\sin x\sin y$；外边界非零，应由外部真解计算。

\section{实验与数据对应}
项目根目录名为\path{Bachelor_Thesis_RaNN}。实验01原始数据保存在\path{results/raw}，其余保存在\path{results/runs}。本文主要子目录为：
\begin{itemize}
\item \path{exp02_interface_20261013}：界面与尺度，150次；
\item \path{exp03_sampling_20261013}：主采样和指标消融，180次；
\item \path{exp04_sensitivity_20261013}：固定比例与单因素，285次；
\item \path{exp06_peaks_20261013}：峰值与unused池，400次；
\item \path{exp07_stability_20261013}：15组谱、90组扰动；
\item \path{exp08_efficiency_20261013}：25次串行成本复测；
\item \path{exp09_reproduction_20261013}：110次初始复现/诊断假设；
\item \path{exp09_literal_20261013}：60次较贴原文字面的初始化追加。
\end{itemize}
每个run保存配置、逐次CSV、均值/样本标准差、输出系数和日志。图可由CSV独立重绘；重复run\_id拒绝覆盖。程序状态ok仅表示正常完成，不代表高精度或理论证明。论文表图由\path{experiments/build_thesis_assets.py}生成，并在\path{thesis/asset_sources.csv}记录来源。

\section{运行及编译}
Python依赖在requirements与lock文件。Windows与Linux解释器路径分别为：
\begin{verbatim}
.venv/Scripts/python.exe
.venv/bin/python
\end{verbatim}
在项目根目录运行测试与实验入口，例如：
\begin{verbatim}
python -m unittest discover -s tests -v
python -m experiments.exp09_reproduction
\end{verbatim}
新run\_id另行指定，不覆盖已经保存的实验结果。

论文推荐XeLaTeX。学校模板尚未提供；有ctexbook时使用该类，否则以book、fontspec和Noto中文字体构成可编译草稿。当前排版及数字顺序参考文献组织只用于审阅，需要按学校规范调整封面、声明、页眉、目录和引用细节。所有实际编译和质量检查结果在项目文档记录，不能把编译成功理解为论文已获学校审核。

\backmatter
% 数字顺序制审阅稿，元数据来自保存的Crossref/arXiv来源；学校细则待确认。
\begin{thebibliography}{99}
\bibitem{dgpreprint} SUN J, DONG S, WANG F. Local randomized neural networks with discontinuous Galerkin methods for partial differential equations[EB/OL]. arXiv:2206.05577v1, 2022[2026-10-13]. \url{https://arxiv.org/abs/2206.05577v1}. 正式版：Journal of Computational and Applied Mathematics, 2024, 445:115830. DOI: \url{https://doi.org/10.1016/j.cam.2024.115830}.

\bibitem{interfacepreprint} LI Y, WANG F. Local Randomized Neural Networks Methods for Interface Problems[EB/OL]. arXiv:2308.03087v1, 2023[2026-10-13]. \url{https://arxiv.org/abs/2308.03087v1}. 正式版题名：Local randomized neural networks with finite difference methods for interface problems, Journal of Computational Physics, 2025, 529:113847. DOI: \url{https://doi.org/10.1016/j.jcp.2025.113847}. 本文算法阅读和复现依据预印本。

\bibitem{ag} DANG H, WANG F, JIANG S. Adaptive-Growth Randomized Neural Networks for PDEs: Algorithms and Numerical Analysis[EB/OL]. arXiv:2408.17225，首次2024，本文读取v3（2025-11-23）[2026-10-13]. \url{https://arxiv.org/abs/2408.17225v3}. DOI: \url{https://doi.org/10.48550/arXiv.2408.17225}.

\bibitem{schwarzpreprint} SHANG Y, HEINLEIN A, MISHRA S, WANG F. Overlapping Schwarz preconditioners for randomized neural networks with domain decomposition[EB/OL]. arXiv:2412.19207v1, 2024[2026-10-13]. \url{https://arxiv.org/abs/2412.19207v1}. 正式版：Computer Methods in Applied Mechanics and Engineering, 2025, 442:118011. DOI: \url{https://doi.org/10.1016/j.cma.2025.118011}.

\bibitem{rar} WU C, ZHU M, TAN Q, KARTHA Y, LU L. A comprehensive study of non-adaptive and residual-based adaptive sampling for physics-informed neural networks[J]. Computer Methods in Applied Mechanics and Engineering, 2023, 403:115671. DOI: \url{https://doi.org/10.1016/j.cma.2022.115671}. 本文读取公开预印本arXiv:2207.10289。

\bibitem{christoffel} LINGHU J, WANG Y. Residual-Christoffel Sampling for Random Feature Collocation of Linear PDEs[EB/OL]. arXiv:2607.13382v1, 2026[2026-10-13]. \url{https://arxiv.org/abs/2607.13382v1}.

\bibitem{pg} SHANG Y, WANG F, SUN J. Randomized neural network with Petrov--Galerkin methods for solving linear and nonlinear partial differential equations[J]. Communications in Nonlinear Science and Numerical Simulation, 2023, 127:107518. DOI: \url{https://doi.org/10.1016/j.cnsns.2023.107518}. 正式书目已核实，本文未获全文，未据此归纳方法或作者结论。

\bibitem{golub} GOLUB G H, VAN LOAN C F. Matrix computations[M]. 4th ed. Johns Hopkins University Press, 2013. ISBN:9781421407944. DOI: \url{https://doi.org/10.56021/9781421407944}. 版本和年份按OpenLibrary第4版记录；本项目未取得教材全文，正文只使用并证明标准线性代数结论。

\bibitem{scipy} VIRTANEN P, GOMMERS R, OLIPHANT T E, et al. SciPy 1.0: fundamental algorithms for scientific computing in Python[J]. Nature Methods, 2020, 17:261--272. DOI: \url{https://doi.org/10.1038/s41592-019-0686-2}.
\end{thebibliography}

\end{document}
